\begin{solapartat}
\punts{0,5} $|\vec{E}| = \frac{\Delta V}{d} = \frac{70.000}{0,02} = 3,5\times 10^{6}\ \mathrm{V/m} = 3,5\times 10^{6}\ \mathrm{N/C}$

\punts{0,25} $|\vec{F}| = eE = 1,602\times 10^{-19}\times 3,5\times 10^{6} = 5,61\times 10^{-13}\ \mathrm{N}$

\medskip
\punts{0,5} Cal dibuixar correctament les línies de camp, perpendiculars a les plaques, dirigides cap el càtode i equidistants atès que el camp és uniforme. Si el sentit és incorrecte, cal restar 0,2 punts; si no es dibuixen equidistants o simplement es representa el camp amb un únic vector enlloc de línies de camp, cal restar 0,2 punts.

Com que la càrrega de l'electró és negativa, el sentit d'$\vec{E}$ és l'oposat del de la força $\vec{F}$. Si no és representa correctament $\vec{F}$, cal restar 0,1 punts.
\figura[0.38]{sol_fig1.png}
\end{solapartat}

\begin{solapartat}
\punts{0,75} Es poden abordar dos plantejaments diferents: un basat en les energies i un altre en les forces.

\destaca{Per energies:} segons el teorema de les forces vives, la variació d'energia cinètica és igual al treball del camp elèctric:
\[ \Delta E_C = W_e = -\Delta U = -q\Delta V = |e|\Delta V = 1,602\times 10^{-19}\times 70.000 = 1,12\times 10^{-14}\ \mathrm{J} \]
\[ \Delta E_C = \frac{1}{2} m v^2 - 0 \Longrightarrow v = \sqrt{\frac{2\cdot 1,12\times 10^{-14}}{9,11\times 10^{-31}}} = 1,57\times 10^{8}\ \mathrm{m/s} \]

\destaca{Per forces:} segons la segona llei de Newton.
\[ a = \frac{F}{m} = \frac{5,61\times 10^{-13}}{9,11\times 10^{-31}} = 6,15\times 10^{17}\ \mathrm{m/s^2} \]
\[ v^2 - 0 = 2\cdot a\cdot d \Rightarrow v = \sqrt{2\cdot a\cdot d} = \sqrt{2\cdot 6,15\times 10^{17}\cdot 0,02} = 1,57\times 10^{8}\ \mathrm{m/s} \]

\punts{0,25} L'energia del fotó és:
\[ E_f = hf = h\,\frac{c}{\lambda} = 6,63\times 10^{-34}\,\frac{3,00\times 10^{8}}{1,00\times 10^{-9}} = 1,99\times 10^{-16}\ \mathrm{J} \]

\punts{0,25} Per conservació de l'energia, podem calcular l'energia cinètica després del xoc:
\begin{align*}
E_{C,i} = E_{C,f} + E_f \Longrightarrow E_{C,f} &= 1,121\times 10^{-14}\ \mathrm{J} - 1,99\times 10^{-16}\ \mathrm{J}\\
&= 1,10\times 10^{-14}\ \mathrm{J}
\end{align*}
I la velocitat de l'electró després del xoc és:
\[ v_f = \sqrt{\frac{2\cdot 1,101\times 10^{-14}}{9,11\times 10^{-31}}} = 1,56\times 10^{8}\ \mathrm{m/s} \]
\end{solapartat}
